% =========================================================
% Completeness Under Equivalent Metrics
% Source: volume-iii/analysis/metric-spaces/notes/metrics/
%
% Purpose: Show via concrete example that equivalent metrics
%          can disagree on completeness.
% Dependencies: def:equivalent-metrics, def:complete-metric-space
% =========================================================

\subsubsection{Completeness is Not a Topological Property}
\label{sec:completeness-not-topological}

Definition~\ref{def:equivalent-metrics} noted that equivalent metrics
induce the same open sets, the same convergent sequences, and the same
compact sets --- but they \emph{may disagree on completeness}.
This subsection makes that concrete.

\begin{remark*}[What ``equivalent'' preserves and what it does not]
Two equivalent metrics $d_1 \sim d_2$ on $X$ agree on:
\begin{itemize}
  \item Which sets are open or closed.
  \item Which sequences converge and to what limits.
  \item Continuity of functions.
  \item Compactness and connectedness.
\end{itemize}
They do \emph{not} necessarily agree on:
\begin{itemize}
  \item Which sequences are Cauchy (Cauchy depends on the specific distances,
        not just on whether distances are small or large).
  \item Whether every Cauchy sequence converges (completeness).
  \item Uniform continuity of functions.
\end{itemize}
Completeness is a \emph{metric} property, not a \emph{topological} property.
\end{remark*}

\begin{example}[Two equivalent metrics on $(0,1)$: one complete, one not]
\label{ex:completeness-metric-dependent}
Let $X = (0,1)$.

\medskip
\noindent\textbf{Metric $d_1$: the inherited absolute-value metric.}
\[
d_1(x,y) = |x - y|.
\]
The sequence $x_n = \tfrac{1}{n}$ satisfies $d_1(x_m, x_n) = \bigl|\tfrac{1}{m} - \tfrac{1}{n}\bigr| \to 0$,
so it is Cauchy in $(X, d_1)$.
But $x_n \to 0 \notin (0,1)$, so it has no limit in $X$.
Therefore $(X, d_1)$ is \textbf{not complete}.

\medskip
\noindent\textbf{Metric $d_2$: the transported metric from $\mathbb{R}$ via $\tan$.}

The map $\varphi : (0,1) \to \mathbb{R}$ defined by
\[
\varphi(x) = \tan\!\left(\pi x - \tfrac{\pi}{2}\right)
\]
is a homeomorphism (a bicontinuous bijection) from $(0,1)$ onto $\mathbb{R}$.
Define
\[
d_2(x, y) = |\varphi(x) - \varphi(y)|.
\]
Since $|\cdot|$ is a metric on $\mathbb{R}$ and $\varphi$ is a bijection,
$d_2$ is a metric on $(0,1)$.

\medskip
A sequence is Cauchy in $(X, d_2)$ if and only if $(\varphi(x_n))$ is Cauchy in
$(\mathbb{R}, |\cdot|)$, hence convergent in $\mathbb{R}$ (since $\mathbb{R}$ is complete).
Since $\varphi$ is a bijection onto $\mathbb{R}$, the limit is $\varphi(x)$ for some
$x \in (0,1)$, and $x_n \to x$ in $(X, d_2)$.
So $(X, d_2)$ is \textbf{complete}.

\medskip
\noindent\textbf{The metrics are equivalent.}

$\varphi$ is a homeomorphism, so $d_1$ and $d_2$ induce the same topology on $(0,1)$:
a set is $d_1$-open if and only if it is $d_2$-open.
They are therefore equivalent metrics (in the topological sense).

\medskip
\noindent\textbf{Summary.}
\[
d_1 \sim d_2
\quad\text{(same topology on } (0,1)\text{)},
\]
but $(X, d_1)$ is incomplete while $(X, d_2)$ is complete.
Changing the metric --- even to an equivalent one --- can create or
destroy completeness.
\end{example}

\begin{remark*}[Intuition: what Cauchy ``means'' in each metric]
Under $d_1$, the sequence $x_n = 1/n$ in $(0,1)$ is Cauchy because the
terms are getting close to each other in the absolute-value sense.
But the sequence is ``running off the edge'' of $(0,1)$ toward $0$.

Under $d_2$, the same sequence \emph{also} has terms getting close to each other
in the $d_2$ sense --- but the map $\varphi$ stretches $(0,1)$ near $0$ and $1$
into the ends of $\mathbb{R}$.
So the sequence $1/n$ corresponds to $\varphi(1/n) = \tan(\pi/n - \pi/2) \to -\infty$,
which is not Cauchy in $(\mathbb{R}, |\cdot|)$ and therefore not Cauchy in $(X, d_2)$.

This is the key insight: $d_2$ sees the approach to the boundary of $(0,1)$
as ``going to infinity,'' whereas $d_1$ does not.
\end{remark*}

\begin{remark*}[Can every incomplete space be completed by re-metrising?]
Yes --- the Completion Theorem (Theorem~\ref{thm:completion}) shows that
every metric space $(X, d)$ can be isometrically embedded in a complete space
$(\widetilde{X}, \widetilde{d})$.
The metric $\widetilde{d}$ on $\widetilde{X}$ restricts to $d$ on $X$, so
it is the \emph{same} metric (not an equivalent one) on the original space.
The completion works by adding the missing limit points, not by changing the metric.
\end{remark*}
